\begin{tikzpicture}[
  node distance=9mm and 8mm,
  block/.style={draw, rounded corners, thick, align=center, minimum width=28mm, minimum height=9mm, fill=black!3},
  side/.style={draw, rounded corners, align=center, minimum width=24mm, minimum height=8mm, fill=black!2},
  arrow/.style={-{Latex[length=2.5mm]}, thick, shorten <=2pt, shorten >=2pt}
]
\node[block] (observe) {Observe\\state};
\node[block, right=of observe] (model) {Model\\boundaries};
\node[block, right=of model] (schedule) {Schedule\\and place};
\node[block, right=of schedule] (update) {Update and\\retain};
\path (observe.north) ++(0,8mm) coordinate (feedbackL);
\path (update.north) ++(0,8mm) coordinate (feedbackR);

\draw[arrow] (observe) -- (model);
\draw[arrow] (model) -- (schedule);
\draw[arrow] (schedule) -- (update);
\draw[arrow] (update.north) -- (feedbackR) -- node[midway, above=2pt, fill=white, inner sep=1pt, font=\small] {telemetry feedback} (feedbackL) -- (observe.north);

\node[side, below=10mm of model] (hardware) {hardware topology};
\node[side, below=10mm of schedule] (service) {latency, energy, quality};
\node[side, below=10mm of update] (memory) {retention and reuse};

\draw[arrow] (hardware.north) -- (model.south);
\draw[arrow] (service.north) -- (schedule.south);
\draw[arrow] (memory.north) -- (update.south);
\end{tikzpicture}
